\section{Conclusões}
\label{conclusoes}

Este trabalho modelou 3.630 alunos como vértices de um grafo de semelhança construído sem o desfecho e comparou duas distâncias e dois métodos em três janelas. Sobre as distâncias, a escolha pesou menos do que H2 previa: a associação produz comunidades mais coesas e as separa melhor no 1º semestre, mas não altera a previsão da GraphSAGE, e as duas distâncias apontam o mesmo grupo de risco na matrícula. Sobre os métodos, a GraphSAGE (V2 e V4) foi a mais precisa das quatro variações em todas as janelas, e a V3 (associação + Leiden), a que melhor descreve perfis. Sobre as janelas, o maior ganho de informação ocorre entre a matrícula e o fim do 1º semestre, o momento mais promissor para intervir, já que a janela seguinte acrescenta pouco e registra parte da evasão em curso.

Os resultados sustentam que, nesta base, a semelhança entre alunos carrega informação sobre a evasão e organiza perfis interpretáveis e estáveis. Não sustentam conclusões causais: mostram quais perfis se associam à evasão, não por que esses alunos evadem. Os dados vêm de uma única instituição, com duas classes de desfecho; a estimativa dos \textit{Enrolled} não tem gabarito e supõe que o padrão se mantém; a GraphSAGE usou configuração padrão; e o grafo inclui todos os alunos, o que corresponde a analisar uma coorte já matriculada. Os próximos passos são testar outros métodos no mesmo grafo (por exemplo, Louvain e GCN), ajustar os hiperparâmetros da GraphSAGE, inserir cada ingressante novo no grafo sem reconstruí-lo e repetir o estudo com dados de outra instituição, como uma universidade brasileira.

Na prática, os perfis podem orientar o acompanhamento de ingressantes mais velhos com pais de escolaridade desconhecida desde a matrícula, e de alunos com poucas aprovações e mensalidade atrasada ao fim do 1º semestre. Entre os 794 alunos ainda matriculados, o perfil da matrícula indica de 284 a 299 evasões esperadas (36\% a 38\%): o Leiden indica quais perfis acompanhar, e a GraphSAGE, por quais alunos começar.

\textbf{Papéis da equipe.} Erick A. Bonifácio (coordenador): cronograma, integração dos artefatos, repositório e aderência às diretrizes. Augusto Juiz Ribeiro (modelador A): delimitação do problema, modelagem formal do grafo, hipóteses e revisão do relatório. Pedro Henrique Mendes da Silva (modelador B): curadoria do conjunto de dados, apoio à modelagem e interpretação dos perfis. Wesley Batista Luz (desenvolvedor A): arquitetura do protótipo e núcleo algorítmico. João Paulo da Costa Santos (desenvolvedor B): implementação, testes e validação quantitativa.
